% Efficient resource allocation for process execution is fundamental in various industries, such as transport, distribution, and manufacturing. We can model these allocation processes as optimization problems. A specific case is the \textit{Just-In-Time Job Shop Scheduling} (JITJSS) problem, where tasks have specific due dates, and completing them outside these deadlines generates earliness or tardiness penalties. The objective is to minimize the sum of these penalties.
% The JITJSS is known to be a strongly NP-Hard problem, meaning that finding an exact solution is computationally infeasible in practice. Therefore, an alternative is to use heuristics, strategies that obtain solutions in a reasonable time, although without optimality guarantees. The challenge in heuristic construction is that the designer must manually test different combinations of algorithmic components to find the most efficient one. This process is tedious, costly, difficult to reproduce, and often inconclusive.
% To mitigate these issues, the scientific community has developed tools and techniques for Automatic Algorithm Configuration (AAC) and Automatic Algorithm Design (AAD). These techniques systematize and automate the mechanical aspects of heuristic development, enabling designers to create more efficient algorithms. The effectiveness of AAD in generating state-of-the-art algorithms is well-documented across various domains. However, these techniques have not yet been applied to the JITJSS problem.
% Therefore, the objective of this work is to apply AAC and AAD techniques to JITJSS, proposing a new algorithmic design space and automatically generating new algorithms for the problem. We expect to obtain solutions that surpass the current state-of-the-art and validate the effectiveness of the constructed components.

% \keywords{Automatic Algorithm Design, Automatic Algorithm Configuration, Scheduling Problems, Heuristics}

Efficient resource allocation for process execution is fundamental in various industries, such as transport, distribution, and manufacturing. We can model these allocation processes as optimization problems. A specific case is the \textit{Just-In-Time Job Shop Scheduling} (JITJSS) problem, where tasks have specific due dates, and completing them outside these deadlines generates earliness or tardiness penalties. The objective is to minimize the sum of these penalties. The JITJSS is known to be a strongly NP-Hard problem, meaning that finding an exact solution is computationally infeasible in practice. Therefore, an alternative is to use heuristics, strategies that obtain solutions in a reasonable time, although without optimality guarantees. The challenge in heuristic construction is that the designer must manually test different combinations of algorithmic components to find the most efficient one. This process is tedious, costly, difficult to reproduce, and often inconclusive. To mitigate these issues, the scientific community has developed tools and techniques for Automatic Algorithm Configuration (AAC) and Automatic Algorithm Design (AAD). These techniques systematize and automate the mechanical aspects of heuristic development, enabling designers to create more efficient algorithms. The effectiveness of AAD in generating state-of-the-art algorithms is well-documented across various domains. However, these techniques have not yet been applied to the JITJSS problem. Therefore, the objective of this work is to apply AAC and AAD techniques to JITJSS, proposing a new algorithmic design space and automatically generating new algorithms for the problem. To this end, a search space parameterized via a context-free grammar was structured. The configuration process converged to an architecture based on Iterated Tabu Search, validating the relevance of the proposed original neighborhood, \texttt{troca\_crit}, which was present in 80\% of the elite solutions. In experimental tests conducted on 72 literature test instances, the generated algorithm matched or surpassed the best known value in approximately 72.2\% of the scenarios. Furthermore, the method demonstrated high temporal efficiency, locating the final solutions in an average of 21.1 seconds. The obtained results validate the effectiveness of the constructed components and endorse automatic design as a robust alternative to manual heuristic tuning for complex scheduling problems.

\keywords{Automatic Algorithm Design, Automatic Algorithm Configuration, Scheduling Problems, Heuristics}